\section{Limitations}
\label{sec:limitations}

The present study is intended as a controlled computational investigation
rather than as a demonstration of a generally applicable adaptive-meshing
method. Several limitations must therefore be considered when interpreting
the results.

\subsection{Fine-grid numerical reference}

The M160 solution is used throughout the adaptive-mesh analysis as a
fine-grid numerical reference. It should not be interpreted as an exact
solution.

Consequently, quantities reported as errors relative to M160 are more
precisely numerical discrepancies with respect to the finest solution used
in this study. Although the uniform-mesh sequence showed strong convergence
toward M160, a remaining discretization error is necessarily present in the
reference itself.

The results therefore support comparisons between refinement strategies
under a common reference, but they do not establish absolute discretization
errors.

\subsection{Spatial and temporal discretization}

The uniform-grid convergence study changed the time step together with the
spatial resolution. The reported M20--M40, M40--M80, and M80--M160
differences therefore do not isolate spatial discretization error
independently from temporal discretization effects.

The observed convergence rates are consistent with strong numerical
convergence of the tested sequence, particularly on the finer meshes, but
they should not be interpreted as a rigorous spatial order-of-accuracy
verification.

A future verification study should independently refine the time step and
the spatial mesh.

\subsection{Lid-driven-cavity corner behavior}

The moving lid meets stationary side walls at the upper corners of the
cavity. This discontinuity in the imposed wall velocity generates strong
local gradients and influences pointwise discrepancy measures.

The largest M40--M160 velocity discrepancy occurs close to the upper-right
corner, and a substantial fraction of the squared discrepancy is
concentrated in a relatively small number of cells near the upper boundary.

This behavior can influence both conventional gradient indicators and
SOM-derived quantities. In particular, the observed correlation between
quantization error and the M40--M160 discrepancy is partly associated with
both quantities responding strongly to extreme upper-corner flow states.

For this reason, maximum pointwise discrepancy is not used as the primary
criterion for comparing adaptive strategies. Volume-weighted mean and RMS
metrics provide more representative global measures for the present study.

\subsection{Single flow configuration}

All adaptive experiments were performed for the two-dimensional
lid-driven-cavity problem at

\begin{equation}
Re=100.
\end{equation}

The conclusions therefore apply only to the tested configuration.

No evidence is presently available to determine whether the relative
behavior of Gradient, SOM-QE, and Hybrid-v1 persists at different Reynolds
numbers, for different geometries, in three-dimensional flows, or for
problems involving heat transfer, turbulence, compressibility, multiphase
phenomena, or other physical processes.

Generalization must therefore be tested rather than assumed.

\subsection{SOM architecture and feature selection}

The adaptive CFD experiments in this investigation use a nominal
\(4\times4\) SOM containing 16 neurons and trained using the selected
velocity and velocity-gradient features. A separate sensitivity study
evaluated \(3\times3\), \(4\times4\), \(5\times5\), and \(6\times6\)
architectures across five random initializations. That analysis quantified
changes in representation error, partition stability, interface density,
and spatial stability of the high-priority interface seeds.

The sensitivity study does not establish an optimal SOM architecture and
does not repeat the complete adaptive CFD calculation for every architecture
and initialization. Consequently, it characterizes robustness of the learned
SOM partition and refinement-selection structure, rather than robustness of
the final adaptive CFD accuracy across all tested SOM configurations.

Similarly, the failure of direct SOM-QE refinement should not be interpreted
as a general failure of self-organizing maps for adaptive CFD. It establishes
only that quantization error from the present SOM configuration did not
provide a superior direct refinement criterion for the tested problem.

Alternative features, SOM dimensions, training procedures, or derived
indicators may produce different behavior. Such alternatives should,
however, be defined without tuning against the numerical reference used for
final evaluation.

\subsection{Hybrid-v1 is not optimized}

Hybrid-v1 uses

\begin{equation}
I_{\mathrm{hyb}}
=
G^{*}(1+0.5T).
\end{equation}

The coefficient \(\alpha=0.5\) was intentionally fixed rather than optimized
against M160. This prevents the reference solution from being used both to
design and evaluate the indicator.

Consequently, Hybrid-v1 should be interpreted as a proof-of-concept test of
whether SOM regime-transition information can modify a conventional
indicator in a controlled manner. It is not presented as an optimized
adaptive-meshing algorithm.

The isolated improvement observed at the 10\% budget is insufficient to
establish superiority over Gradient.

\subsection{Refinement-budget range}

Only four refinement budgets were investigated:

\begin{equation}
5\%,\qquad10\%,\qquad15\%,\qquad20\%.
\end{equation}

These experiments are sufficient to reveal important differences in
monotonicity and relative behavior, but they do not fully characterize the
accuracy--cost response of each method.

In particular, the present results should not be extrapolated to very small
or very large refinement fractions without additional simulations.

\subsection{Computational-cost measurements}

The reported execution times correspond to individual CFD solves and were
obtained from single runs. They are therefore susceptible to normal
variability in workstation load and operating-system scheduling.

Furthermore, the timings do not include the complete adaptive workflow,
which may contain:

\begin{itemize}
    \item the pilot CFD calculation;
    \item feature extraction;
    \item SOM training;
    \item indicator calculation;
    \item cell selection;
    \item mesh refinement;
    \item field transfer and post-processing.
\end{itemize}

No claim of computational-speed superiority is therefore made.

A rigorous efficiency study should use repeated timing measurements and
compare total workflow cost for a specified accuracy or quantity-of-interest
target.

\subsection{Refinement strategy}

The current framework performs one refinement operation based on information
from a pilot solution. It does not implement repeated adaptive cycles in
which the solution is recomputed, the indicator reevaluated, and the mesh
adapted iteratively.

The results therefore evaluate one-shot refinement decisions rather than a
fully dynamic adaptive-mesh algorithm.

Multiple adaptation cycles could alter the relative behavior of the tested
indicators and constitute a separate research question.


\subsection{Cell-matched second-stage experiments}

The SOM-interface experiments were deliberately evaluated against
cell-matched Gradient controls because graph expansion produces selected-cell
counts that are not equal to the nominal percentage used in the original
adaptive matrix. The \(1N\) and \(2N\) cases therefore should not be combined
directly with the original 5--20\% budget curves as if they were additional
points on the same budget sequence.

The cell-matched design controls final mesh size, but it does not eliminate
all possible confounding factors. The SOM-interface method uses fixed seeds
and graph expansion, whereas Gradient directly ranks individual cells.
Differences in spatial connectivity and refinement topology are therefore
part of the method being tested.

\subsection{Scope of the conclusions}

The strongest conclusion supported by the present evidence is deliberately
limited.

For the tested \(Re=100\) lid-driven cavity, the conventional
velocity-gradient indicator produced the most consistently accurate
solutions in the original 5--20\% SOM-QE/Hybrid-v1 refinement matrix.
Direct SOM quantization-error refinement was inferior and non-monotonic,
while Hybrid-v1 did not demonstrate a systematic advantage.

A subsequent, separately formulated SOM-interface criterion produced lower
volume-weighted mean and RMS discrepancies than cell-matched Gradient
controls at 1834 and 1942 cells. These experiments are encouraging but
remain limited to one geometry, one Reynolds number, one SOM configuration,
and two graph-expansion choices. They therefore do not establish general
superiority.

The present evidence supports the narrower conclusion that SOM-derived
flow-state interfaces can provide information that is complementary to
velocity-gradient magnitude. Whether this information can be converted into
a robust and general predictor of refinement utility remains unresolved.

